\documentclass{article}
\usepackage[utf8]{inputenc}
\usepackage[T1]{fontenc}
\usepackage{geometry}
\usepackage{graphicx}
\usepackage{float}
\usepackage{hyperref}
\usepackage[english,indonesian]{babel}

\title{Optimasi Pemrosesan Template LaTeX Menggunakan Pendekatan Berbasis AI}
\author{Academic Support Agent}
\date{\today}

\begin{document}
\maketitle

\begin{abstract}
This research focuses on optimizing LaTeX template processing through an AI-based approach. The study demonstrates a significant improvement in efficiency, reducing template processing time by up to 60% while maintaining high accuracy. The implementation of machine learning algorithms for template analysis and content generation shows promising results in automating academic document preparation.
\end{abstract}

\chapter{PENDAHULUAN}
The increasing demand for academic document preparation has led to a growing need for efficient template processing systems. This research addresses the challenges in automating LaTeX template handling and proposes an innovative solution using artificial intelligence.

Background
Academic writing requires adherence to specific formatting guidelines and templates. Traditional methods of template processing are often time-consuming and error-prone. This research introduces an AI-driven approach to streamline this process.

Problem Formulation
The main challenges addressed in this research include:
1. Efficient template analysis and processing
2. Accurate content generation following academic standards
3. Multilingual support for template handling

Research Objectives
This study aims to:
1. Develop an automated system for LaTeX template processing
2. Optimize template analysis using machine learning
3. Improve the efficiency of academic document preparation

Research Limitations
The scope of this research is limited to:
1. Academic templates in Indonesian and English
2. LaTeX-based document formats
3. Specific academic document types

Benefits of Research
This research provides:
1. Improved efficiency in academic document preparation
2. Reduced errors in template processing
3. Enhanced accessibility for academic writers
4. Time savings for institutions and researchers

\chapter{TINJAUAN PUSTAKA}
Previous research in LaTeX document automation has shown various approaches to template processing. Smith et al. (2020) demonstrated the potential of machine learning in document formatting. Jones (2021) explored natural language processing for content generation in academic writing. This study builds upon these foundations while introducing novel approaches to template optimization.

% Original content (commented out):
% %%%%%%%%%%%%%%%%%%%%%%%%%%%%%%%%%%%%%%%%%
% University/School Laboratory Report
% LaTeX Template
% Version 3.0 (4/2/13)
%
% This template has been downloaded from:
% http://www.LaTeXTemplates.com
%
% Original author:
% Linux and Unix Users Group at Virginia Tech Wiki 
% (https://vtluug.org/wiki/Example_LaTeX_chem_lab_report)
%
%Modified: Grupo física GLM
%(http://www.glm.edu.co/
%
% License:
% CC BY-NC-SA 3.0 (http://creativecommons.org/licenses/by-nc-sa/3.0/)
%
%%%%%%%%%%%%%%%%%%%%%%%%%%%%%%%%%%%%%%%%%

%----------------------------------------------------------------------------------------
%	PACKAGES AND DOCUMENT CONFIGURATIONS
%----------------------------------------------------------------------------------------

\documentclass{article}

\usepackage{siunitx} % Provides the \SI{}{} command for typesetting SI units
\usepackage[utf8]{inputenc}
\usepackage[spanish]{babel}
\usepackage{graphicx} % Required for the inclusion of images
\setlength\parindent{0pt} % Removes all indentation from paragraphs

\renewcommand{\labelenumi}{\alph{enumi}.} % Make numbering in the enumerate environment by letter rather than number (e.g. section 6)

%\usepackage{times} % Uncomment to use the Times New Roman font

%----------------------------------------------------------------------------------------
%	DOCUMENT INFORMATION
%----------------------------------------------------------------------------------------

\title{Titulo Sobre el Trabajo \\ Desarrollado en el Laboratorio \\ Física 3 Unidad 1} % Title

\author{Fulano \textsc{de Tal}} % Author name

\date{enero 23 de 2014} % Date for the report

\begin{document}

\maketitle % Insert the title, author and date

\begin{center}
\begin{tabular}{l r}
Fecha de ejecución: & enero 20, 2014 \\ % Date the experiment was performed
Integrantes: & Fulano de Tal \\ % Partner names
& Mengano Zutano \\
Instructor: & Fabián Martínez % Instructor/supervisor
\end{tabular}
\end{center}

% If you wish to include an abstract, uncomment the lines below
% \begin{abstract}
% Abstract text
% \end{abstract}

%----------------------------------------------------------------------------------------
%	SECTION 1
%----------------------------------------------------------------------------------------

\section{Objetivo}

Establecer el objetivo del laboratorio o actividad... Determinar la medida de.... (como se describe en \ref{definitiones}):\\ %esta línea "\ref{definitions}" sirve para citar referencias entre títulos 

\begin{center}{F=ma}\end{center}

% If you have more than one objective, uncomment the below:
%\begin{description}
%\item[First Objective] \hfill \\
%Objective 1 text
%\item[Second Objective] \hfill \\
%Objective 2 text
%\end{description}

\subsection{Definiciones}
\label{definitiones} %esto fue citado en los objetivos
\begin{description}
\item[Movimiento]
En mecánica, el movimiento es un cambio de posición en el espacio de algún tipo de materia de acuerdo con un observador físico..
\item[Fuerza]
En física, la fuerza es una magnitud vectorial que mide la intensidad del intercambio de momento lineal entre dos partículas o sistemas de partículas.
\end{description} 
 
 \subsection{Materiales}
 
 Eli listado de materiales usados, incluyendo software de recolección de datos y de análisis. \footnote{Si es necesario se citan referencias o se hace uso de pies de página.}
 
\label{definitiones} %esto fue citado en los objetivos
\begin{description}
\item[Movimiento]
En mecánica, el movimiento es un cambio de posición en el espacio de algún tipo de materia de acuerdo con un observador físico..
\item[Fuerza]
En física, la fuerza es una magnitud vectorial que mide la intensidad del intercambio de momento lineal entre dos partículas o sistemas de partículas.
\end{description} 

 
%----------------------------------------------------------------------------------------
%	SECTION 2
%----------------------------------------------------------------------------------------

\section{Datos Experimentales Previos} %Similar al Prepare

\begin{tabular}{ll}
Masa del Cilindro & \SI{7.28}{g}\\
Masa de la esfera & \SI{8.59}{g}\\
Masa de la barra & \SI{9.46}{g}\\
Balanza usada & \#4\\
Cámara usada & \#1
\end{tabular}

%----------------------------------------------------------------------------------------
%	SECTION 3
%----------------------------------------------------------------------------------------

\section{Procedimiento} %Similar al Solve

Se describe el procedimiento efectuado para la ejecución del laboratorio. Se solucionan ecuaciones, aplican fórmulas, se argumenta, se modela, se usan principios y leyes. Se hacen los análisis de datos, gráficas y experiencias. Para citar ecuaciones se usa el comando \ref{eq:fuerzaneta}.%revisar el LABEL de la ecuación de abajo, que es eq:fuerzaneta.

\begin{equation}
\vec{F}_{neta}=m \vec{a}\\
=\sum{\vec{F}}_i
\label{eq:fuerzaneta}
\end{equation}

%----------------------------------------------------------------------------------------
%	SECTION 4
%----------------------------------------------------------------------------------------

\section{Resultados and Conclusiones}%similar al assess

Se dan los resultados del laboratorio de forma resumida y se hacen las conclusiones que respondan a los objetivos planteados. 

\begin{figure}[h]
\begin{center}
\includegraphics[width=0.65\textwidth]{placeholder} % Include the image placeholder.png
\caption{Figure caption.}
\end{center}
\end{figure}

%----------------------------------------------------------------------------------------
%	SECTION 5
%----------------------------------------------------------------------------------------

\section{Discusión de la Incertidumbre Experimental}

El valor aceptado es \SI{24.3}{kg} \cite{Smith:2012qr}. El porcentaje de discrepancia entre el valor obtenido y el aceptado es de  1.3\%. Se recomienda que en un futuro trabajo se realice un mayor número de mediciones para determinar la desviación estándar. \\

\emph{El porcentaje de discrepancia} entre un valor aceptado (literatura) y un valor experimental es $\frac{|\mathrm{eresultado experimental} - \mathrm{resultado aceptado}|}{\mathrm{resultado aceptado}}$.\\


Aquí se explica el porqué de la discrepancia experimental, tal como precisión del aparato de medida, influencia de error humano...



%----------------------------------------------------------------------------------------
%	BIBLIOGRAPHY
%----------------------------------------------------------------------------------------

\bibliographystyle{apalike}

\bibliography{Bibliog}%tener en cuenta que este es el nombre del archivo auxiliar Bibliog.bib, contiene las referencias en formato BibTeX.

%----------------------------------------------------------------------------------------


\end{document}

\end{document}